\begin{theorem}[failure of the printed endpoint implication]
\label{thm:nodal-test}
Let $A=\operatorname{pr}_1^*A_\lambda$ on
$S=C_0\times\P^1$, with $\lambda$ as in Lemma~\ref{lem:gluing-bundle}.
Then
\[
 H^0(S,N^*_{S/\mathcal X}\otimes A)=0,
 \qquad
 H^0(S,\widetilde\Omega^1_S\otimes A)=0,
 \tag{4.5}
\]
but
\[
 H^0(S,\Omega^1_{\mathcal X}|_S\otimes A)\ne0.
 \tag{4.6}
\]
In particular, the two endpoint vanishings in (4.5) do not imply the
vanishing of (4.6), even for a proper morphism over a disc with smooth total
space and an irreducible reduced normal-crossing (self-crossing) central
fibre.
\end{theorem}

\begin{proof}
Because $S$ is the principal fibre $(t=0)$, its conormal line is trivial:
\(N^*_{S/\mathcal X}\simeq\cO_S\).  The projection
\(\operatorname{pr}_1:S\to C_0\) has
\(\operatorname{pr}_{1*}\cO_S\simeq\cO_{C_0}\), so (4.4) gives
\[
 H^0(S,A)\simeq H^0(C_0,A_\lambda)=0.
 \tag{4.7}
\]
This proves the first vanishing in (4.5).

We next check the second.  At a normal-crossing (self-crossing) point with local equation
\(uv=0\), the pullback map
\(\Omega^1_S\to\eta_*\Omega^1_{\widetilde S}\) has kernel exactly
\(\mathscr T_S\).  Indeed, in
\(R=\C\{u,v,w\}/(uv)\) the relation is
\(v\,du+u\,dv=0\).  If
\(a\,du+b\,dv+c\,dw\) restricts to zero on both branches, then
\(a\in(v)\), \(b\in(u)\), and
\(c\in(u)\cap(v)=(uv)=0\).  Modulo the relation the class is therefore a
multiple of $\tau=v\,du=-u\,dv$.  This class is nonzero: if
\(v\,du=h(v\,du+u\,dv)\), then coefficient comparison gives
  \((1-h)v=0\) and $hu=0$.  Since
\(\operatorname{Ann}_R(v)=(u)\) and \(\operatorname{Ann}_R(u)=(v)\),
\(1-h\in(u)\) and \(h\in(v)\), which would imply
\(1=(1-h)+h\in(u)+(v)\), impossible because $(u)+(v)$ is proper.
Moreover $(u+v)\tau=0$.  Since $u+v$ is a non-zero-divisor in
$R$ (the associated primes are $(u)$ and $(v)$), $\tau$ is genuine
torsion.
Conversely, each summand of the normalization module is a module over one of
the domains $R/(u)$ and $R/(v)$.  A non-zero-divisor of the reduced ring $R$
avoids both minimal primes and therefore acts injectively on both summands;
the target is torsion-free.  Thus the kernel of the normalization map is
exactly the torsion subsheaf, and there is an injection
\[
 \widetilde\Omega^1_S\hookrightarrow\eta_*\Omega^1_{\P^1\times\P^1}.
 \tag{4.8}
\]
Since $\eta^*A\simeq\cO_{\P^1\times\P^1}$, the projection formula and (4.8)
give
\[
 H^0(S,\widetilde\Omega^1_S\otimes A)
 \hookrightarrow
 H^0(\P^1\times\P^1,\Omega^1_{\P^1\times\P^1})=0.
 \tag{4.9}
\]
The last equality follows from
\(\Omega^1_{\P^1\times\P^1}\simeq
\operatorname{pr}_1^*\Omega^1_{\P^1}\oplus
\operatorname{pr}_2^*\Omega^1_{\P^1}\) and
\(H^0(\P^1,\Omega^1_{\P^1})=0\).

Finally, along the conductor $D$ the pullback of $dg|_S$ is zero.  Locally
$g=h\,uv$ with $h$ a unit, so after restriction to $S$ one has
$i^*(dg)=h(v\,du+u\,dv)$; both coefficients vanish on the preimage of $D$.
Proposition~\ref{prop:conductor-differential}
therefore gives a nonzero section in (4.6).  This proves the theorem.
\end{proof}

\begin{remark}[what the theorem does and does not refute]
\label{rem:scope}
The bundle $A$ is non-torsion on $S$: restriction to
\(C_0\times\{r\}\) would otherwise make $A_\lambda$ torsion.  Thus the
example meets the non-torsion condition attached to the p.~692 reduction.
In fact Lemma~\ref{lem:extending-gluing} gives an extending choice of $A$ on
the displayed smooth total space.  If $\widetilde L:=\omega_{\mathcal X}\otimes
\widetilde A^{-1}$, then
$(\widetilde L^\vee\otimes\omega_{\mathcal X})|_S\simeq A$, so the example
also has the Serre-dual twist appearing in the CDP argument.  This family
nevertheless does not have
the global topological and algebraic-dimension hypotheses of CDP20
Theorem~2.2.  The conclusion is exactly the formal one stated in
Theorem~\ref{thm:nodal-test}: the two displayed endpoint vanishings do not
control the torsion-supported lifting obstruction $\delta_A$.
\end{remark}

\section{The corrected reading of the CDP step}

The published paper states at p.~692 that, for an irreducible reduced fibre,
the conormal sequence makes it sufficient (when the restriction of the line
bundle is non-torsion) to prove the quotient vanishing (1.1)
\cite{CDP2020}.  Propositions~\ref{prop:filtration} and
\ref{prop:criterion} give the exact replacement:
\[
\begin{gathered}
 H^0(S,\widetilde\Omega^1_S\otimes A)=0,
 \quad H^0(S,N^*_{S/X}\otimes A)=0,\\[-2pt]
 \Longrightarrow\quad
 H^0(S,\Omega^1_X|_S\otimes A)=\ker\delta_A.
\end{gathered}
\tag{5.1}
\]
The missing assertion is therefore not another endpoint vanishing but the
injectivity of the explicitly defined map $\delta_A$.  For a reducible
reduced fibre the same formula must be applied to the whole Cartier fibre,
not separately to its components.  For a nonreduced fibre there are further
nilpotent layers, and (5.1) is not a complete analysis of that case.

The test family above is a counterexample to the unrestricted three-sheaf
implication for an arbitrary fibre line bundle $A$, not to CDP20
Theorem~2.2 or to its singular-fibre proposition under its full hypotheses.
Even though $A$ extends over this local total space, the family is not a
compact CDP threefold; whether the global hypotheses force a different
argument to make $\delta_A$ injective remains a separate question.

\subsection{A derived global criterion that survives the correction}

For completeness, we record a derived reformulation of the part of the CDP
strategy that does not use the defective endpoint reduction
\cite[Prop.~2.4]{CDP2020}.  The following generic-vanishing lemma is
elementary but useful in isolating the exact extra hypothesis; no priority is
claimed for this reformulation.

